\SourceLeaf{470}
\par\Needspace{9\baselineskip}\bigskip
\pdfbookmark[3]{Darstellung, 20}{pallas-II-article20}
\begin{center}20.\end{center}
\begingroup
\setlength{\lineskip}{4pt}\setlength{\lineskiplimit}{4pt}
\noindent\hspace*{\parindent}Die veränderlichen Werte der Elemente, die man mit dieser Methode erhält, sind bis auf Größen von der Ordnung des Quadrats der störenden Massen richtig. Wünscht man eine noch größere Genauigkeit, so braucht man dieselbe Rechnung nur mit den eben gefundenen veränderlichen Elementen zu wiederholen. So erhält man Ergebnisse, die bis auf Größen dritter Ordnung richtig sind. Es wäre überflüssig, die Rechnung noch einmal zu wiederholen. Bei der zweiten Rechnung darf man sogar annehmen. dass die Elemente eine Zeit lang konstant bleiben, das heißt, bei der Berechnung der momentanen Änderungen ein und dasselbe System von Elementen für mehrere Epochen vor oder nach derjenigen verwenden, der es entspricht. Im Allgemeinen wechselt man das System erst, wenn der Unterschied zwischen ihm und dem System, das an seine Stelle treten müsste, hinreichend merklich geworden ist. Je kleiner also die momentanen Änderungen der Elemente sind, desto größer ist die Zahl der Epochen, für die man dasselbe System von Elementen verwenden darf.
\par Für die Praxis ist es etwas bequemer, die Werte von \(\dfrac{\mathrm{d}i}{\mathrm{d}t},\ \dfrac{\mathrm{d}\AscendingNode}{\mathrm{d}t},\ \dfrac{\mathrm{d}\varphi}{\mathrm{d}t},\ \dfrac{\mathrm{d}\tilde\omega}{\mathrm{d}t},\ \dfrac{\mathrm{d}\epsilon}{\mathrm{d}t}\) mit \(\theta\) und die Werte von \(\dfrac{\mathrm{d}n}{\mathrm{d}t}\) mit \(\theta\theta\) zu multiplizieren, bevor man die Differenzen- und Summenreihen bildet. Es ist überflüssig zu erwähnen, dass man diese Größen zunächst in Sekunden ausdrücken sollte. Die zweckmäßigste Länge des Intervalls \(\theta\) hängt teils von der Geschwindigkeit der Bewegung des störenden und des gestörten Planeten ab, teils vom angestrebten Grad der Genauigkeit. Als der Verfasser mit dieser Methode die durch Jupiter verursachten Störungen der Pallas in den Jahren \mbox{1803 . . . . 1816} berechnete, wählte er ein Intervall von 50 Tagen. Ein verständiger Rechner hat zu beurteilen, ob man dieses Intervall unter anderen Umständen verkürzen oder verlängern sollte. Wir müssen hier mehrere weitere praktische Bemerkungen übergehen, die sich jedem geübten Rechner bei ähnlichen Arbeiten von selbst ergeben werden.
\par\endgroup
\begin{center}\rule{0.18\textwidth}{0.35pt}\end{center}
